\section{Introduction}

Election forensic analytics ask whether reported results contain unusual
patterns: precinct swings that do not fit their history, turnout or vote-share
residuals far from comparable units, batch and scanner effects, leading- or
last-digit distributions that look improbable, neighbors that diverge
unexpectedly, abrupt change points across time or batches, and audit
discrepancies that cluster by machine or source. These analytics are useful for
prioritizing where humans should look. They do not, by themselves, certify an
outcome or prove that anything went wrong.

That distinction is the reason for a separate W-series. The V-series verifies
canvass and audit claims with explicit evidence contracts: a passing V-series
transcript means a declared count or audit claim was replayed from public
inputs. A W-series report makes no such claim. It produces an investigation
queue --- a ranked set of units, batches, scanners, or contests that merit more
review --- and it must say so in plain language on its face.

This paper leads with that boundary and then builds the shared scaffolding the
rest of the track depends on. It surveys the analytic families that clean RCOUNT
inputs can feed, specifies a single forensic report contract every method must
satisfy, and lays out the fixture pattern that lets a report be reproduced from
hash-bound inputs. It is a framework paper. The analytic methods it names are
established in the forensics literature; the contribution here is the
RCOUNT-grounded report contract and the strict separation that keeps anomaly
scores from being read as certification evidence.
